%%%%%%%%%%%%%%%%%%%%%%%%%%%%%%%%%%%%%%%%%%%%%%%%%%%%%%%%%%%%%%%%%%%%%%
% Lớp 10 | Phương pháp tọa độ trong mặt phẳng (𝑂𝑥𝑦) | Tọa độ của véctơ
% Dạng 0H9?1-2: Phép toán véctơ (tổng, hiệu, tích với số ) trong Oxy
% Số câu: 11  (sinh tự động bởi tach_cau.py; từ nay thầy sửa trực tiếp file này)
%%%%%%%%%%%%%%%%%%%%%%%%%%%%%%%%%%%%%%%%%%%%%%%%%%%%%%%%%%%%%%%%%%%%%%

%%%%%%%%%%%%%%%% TRẮC NGHIỆM 4 PHƯƠNG ÁN %%%%%%%%%%%%%%%%

% nguồn: MaHoaNH_TracNghiem_Lop10
\begin{ex}%[0H9N1-2]%[ma:de85301c]
	Cho hai vectơ $\overrightarrow{a}=(1;2)$; $\overrightarrow{b}=(3;4)$. Tọa độ $\overrightarrow{c}=4\overrightarrow{a}-\overrightarrow{b}$  bằng
	 \choice
		 {\True $(1;4)$}
		 {$(-1;4)$}
		{$(4;1)$} 
		 {$(-1;-4)$}
	\loigiai{ 
	Ta có $4\overrightarrow{a}=(4;8)$ nên $\overrightarrow{c}=4\overrightarrow{a}-\overrightarrow{b}=(4-3;8-4)=(1;4)$.} 
\end{ex}

% nguồn: MaHoaNH_TracNghiem_Lop10
\begin{ex}%[0H9N1-2]%[ma:67497f3e]
	Trong mặt phẳng tọa  độ $Oxy$, cho vectơ $\overrightarrow{n}=(3;-4)$. Tính vectơ $\overrightarrow{m}=-2\overrightarrow{n}$.
	 \choice
		 {$\overrightarrow{m}=(6;8)$}
		{\True $\overrightarrow{m}=(-6;8)$} 
		{$\overrightarrow{m}=(6;-8)$} 
		 {$\overrightarrow{m}=(-6;-8)$}
	 \loigiai{
		Ta có  $\overrightarrow{m}=-2\overrightarrow{n}=(-6;8)$.
	} 
\end{ex}

% nguồn: MaHoaNH_TracNghiem_Lop10
\begin{ex}%[0H9N1-2]%[ma:9c93ff14]
	Trong mặt phẳng tọa độ $Oxy$, cho hai vectơ $\overrightarrow{a}=\left(0;3\right)$ và  $\overrightarrow{b}=\left(4;-6\right)$. Tìm tọa độ của vectơ $\overrightarrow{c}=\overrightarrow{a}+2\overrightarrow{b}$.
	\choice 
		{$\overrightarrow{c}=\left(8;15\right)$} 
		 {$\overrightarrow{c}=\left(4;-3\right)$}
		{\True  $\overrightarrow{c}=\left(8;-9\right)$}
		{$\overrightarrow{c}=\left(8;-3\right)$} 
	\loigiai{ Gọi $\overrightarrow{c}=(x;y)$, ta  có
		$ \overrightarrow{c}=\overrightarrow{a}+2\overrightarrow{b} \Leftrightarrow \heva{&x=0+2\cdot 4=8\\&y=3+2\cdot (-6)=9}\Rightarrow \overrightarrow{c}=\left(8;-9\right) $. 
	 }
\end{ex}

% nguồn: MaHoaNH_TracNghiem_Lop10
\begin{ex}%[0H9N1-2]%[ma:880883a4]
	Trong mặt phẳng toạ độ $Oxy$, cho điểm  $A(1;-4)$, điểm $B(2;-1)$. Toạ độ vectơ $\overrightarrow{AB}$ là
	 \choice
		 {$\overrightarrow{AB}=(-1;-3)$}
		 {$\overrightarrow{AB}=(3;-5)$}
		{\True $\overrightarrow{AB}=(1;  3)$}
		{$\overrightarrow{AB}=(1;-3)$} 
	 \loigiai{
		Ta có  $\overrightarrow{AB}=\left(x_B-x_A; y_B-y_A\right)=(2-1;-1--4)=(1; 3)$.
	} 
\end{ex}

% nguồn: MaHoaNH_TracNghiem_Lop10
\begin{ex}%[0H9N1-2]%[ma:429e8229]
	Trong mặt phẳng toạ độ $Oxy$, cho $\overrightarrow{u}=(5;-4)$. Biểu diễn của véc-tơ $\overrightarrow{u}$ qua véc-tơ $\overrightarrow{i}$  và $\overrightarrow{j}$ là
	 \choice
		{$\overrightarrow{u}=-5\overrightarrow{i}-4\overrightarrow{j}$} 
		{\True $\overrightarrow{u}=5\overrightarrow{i}-4\overrightarrow{j}$} 
		 {$\overrightarrow{u}=-4\overrightarrow{i}-5\overrightarrow{j}$}
		{$\overrightarrow{u}=-5\overrightarrow{i}+4\overrightarrow{j}$} 
	\loigiai{ 
		Ta có $\overrightarrow{u}=(5;-4)$ nên $\overrightarrow{u}=5\overrightarrow{i}-4\overrightarrow{j}$. 
	} 
\end{ex}

% nguồn: MaHoaNH_TracNghiem_Lop10
\begin{ex}%[0H9N1-2]%[ma:5199186b]
	Trong hệ tọa độ $Oxy$, cho $A(2; 5)$, $B(1; 1)$, $C(3; 3)$. Tìm tọa độ điểm $E$ sao cho $\overrightarrow{AE}=3  \overrightarrow{AB}-2 \overrightarrow{AC}$.
	\choice 
		 {$(3;-3)$}
		 {$(-3; 3)$}
		{\True  $(-3;-3)$}
		 {$(-2;-3)$}
	\loigiai{Gọi $E(x; y)$. 
		Ta có $\overrightarrow{AE}=3 \overrightarrow{AB}-2 \overrightarrow{AC} \Leftrightarrow \overrightarrow{AE}-\overrightarrow{AB}=2(\overrightarrow{AB}-\overrightarrow{AC}) \Leftrightarrow \overrightarrow{BE}=2  \overrightarrow{CB}$.\\
		Suy ra $(x-1;  y-1)=2(-2;-2) \Leftrightarrow\heva{x-1=-4 \\ y-1=-4} \Leftrightarrow\heva{x=-3 \\ y=-3.}$\\
	Vậy  $E(-3;-3)$. }
\end{ex}

% nguồn: MaHoaNH_TracNghiem_Lop10
\begin{ex}%[0H9N1-2]%[ma:07a98557]
	Trong mặt phẳng tọa độ $Oxy$,  cho $ \overrightarrow{a}=(1;-2) $ và $ \overrightarrow{b}=(3;4)$. Véc-tơ $ \overrightarrow{m}=2\overrightarrow{a}+3\overrightarrow{b} $ có tọa độ là
	\choice 
		{$  \overrightarrow{m}=(10;12) $}
		{$ \overrightarrow{m}=(11;16)  $}
		{\True $ \overrightarrow{m}=(11;8)  $}
		{$ \overrightarrow{m}=(11;2) $} 
	 \loigiai{
		Ta có $ 2\overrightarrow{a}=(2;-4) $,  $ 3\overrightarrow{b}=(9;12) $.
		Khi đó $ \overrightarrow{m}=2\overrightarrow{a}+3\overrightarrow{b}=(11;8)  $.
	 }
\end{ex}

% nguồn: MaHoaNH_TracNghiem_Lop10
\begin{ex}%[0H9H1-2]%[ma:6e07fe96]
	Trên hệ tọa độ $Oxy$, cho hai véc-tơ $\overrightarrow{a}=(2;-1)$, $\overrightarrow{b}=(3;-2)$. Tọa độ của  véc-tơ $\overrightarrow{b}-\overrightarrow{a}$ là
	\choice 
		{$(1;-3)$} 
		{$(-1;-3)$} 
		{\True $(1;-1)$} 
		{$(-1;1)$} 
	 \loigiai{
		Ta có $\overrightarrow{b}-\overrightarrow{a}=(3-2;-2+1)=(1;-1)$. 
	} 
\end{ex}

% nguồn: MaHoaNH_TracNghiem_Lop10
\begin{ex}%[0H9H1-2]%[ma:7232406c]
	Cho $\vec{a}=(1; 5)$,  $\vec{b}=(-2; 1)$. Tính $\vec{c}=3\vec{a}+2\vec{b}$.
	\choice 
		{$\vec{c}=(7;  13)$}
		 {$\vec{c}=(1; 17)$}
		{\True $\vec{c}=(-1;  17)$}
		{$\vec{c}=(1;  16)$}
	\loigiai{ 
		 Ta có $\heva{&\vec{a}=(1; 5)\\&\vec{b}=(-2; 1)} \Rightarrow \heva{&3\vec{a}=(3; 15)\\&2\vec{b}=(-4; 2)} \Rightarrow \vec{c}=3\vec{a}+2\vec{b}=(-1; 17)$.
	} 
\end{ex}


%%%%%%%%%%%%%%%% ĐÚNG/SAI %%%%%%%%%%%%%%%%

% nguồn: MaHoaNH_DungSai_Lop10
\begin{ex}%[0H9N1-2]%[ma:3b505586]
	Trên hệ trục tọa độ $Oxy$, cho các  điểm $A(1;2)$, $B(-2;3)$, $C(4;-7)$. Xét tính đúng sai của các mệnh đề sau
	\choiceTF 
		{Trung điểm của  $AB$ có tọa độ là $E(-1;5)$}
		 {\True Ba điểm $A,B,C$ không thẳng hàng}
		{\True Tứ giác $ABCD$ là  hình bình hành với $D(7;-8)$}
		{\True  ${AC}=3\sqrt{10}$}
	 \loigiai{
		\begin{itemchoice}
			\itemch Trung điểm của  $AB$ là $E\left(-\dfrac{1}{2};\dfrac{5}{2} \right) $.
			\itemch  Ta có $\overrightarrow{AB}=(-3;1)$, $\overrightarrow{AC}=(3;-9)$, suy ra hai véc-tơ $\overrightarrow{AB}$, $\overrightarrow{AC}$ không cùng phương. Vì vậy $A,B,C$ không thẳng hàng.
			\itemch $ABCD$ là hình bình hành khi và chỉ khi $\overrightarrow{AB}=\overrightarrow{DC}$. Khi  đó, ta có $D(7;-8)$.
			\itemch Ta có  ${AC}=\sqrt{3^2+9^2}=3\sqrt{10}$.
		\end{itemchoice}
	 }
\end{ex}

% nguồn: MaHoaNH_DungSai_Lop10
\begin{ex}%[0H9T1-2]%[ma:b53eae71]
	Trên hệ trục tọa  độ $Oxy$, cho tam giác $ABC$ với $A(-2;1)$, $B(2;3)$, $C(1;7)$. Xét tính đúng sai của các mệnh đề sau
	\choiceTF 
		{\True Trung điểm của  $AC$ có tọa độ là $F\left(-\dfrac{1}{2};4 \right) $}
		{ Độ dài trung tuyến  ứng với đỉnh $A$ trong tam giác $ABC$ là $\sqrt{85}$}
		{  Tọa độ trực tâm tam giác $ABC$ là $H\left(\dfrac{7}{3};\dfrac{10}{3} \right)$ }
		{\True Trọng tâm tam giác $ABC$ có tọa độ là  $G\left(\dfrac{1}{3};\dfrac{11}{3} \right)$ }
	\loigiai{ 
		\begin{itemchoice}
			\itemch Trung điểm của $AC$ có tọa  độ là $F\left(-\dfrac{1}{2};4 \right) $.
			\itemch Tọa độ trung điểm $M$ của $BC$ là $M\left(\dfrac{3}{2};5 \right)$. Khi đó, độ dài $AM$ là  $\dfrac{\sqrt{85}}{2}$.
			\itemch  Đặt $H(a;b)$ trực tâm của tam giác $ABC$. Ta có $\heva{&\overrightarrow{AH}\cdot \overrightarrow{BC}=0\\&\overrightarrow{BH}\cdot \overrightarrow{AC}=0}\Leftrightarrow \heva{&a=\dfrac{10}{3}\\&b=\dfrac{7}{3}.}$
			\itemch  Trọng tâm tam giác $ABC$ có tọa độ là $G\left(\dfrac{1}{3};\dfrac{11}{3} \right)$.
		\end{itemchoice}
	 }
\end{ex}
